\uuid{OcTp}
\titre{Distance}
\auteur{bodin ; hanani ; mzari}
\organisation{exo7}
\licence{CC BY-NC-SA 4.0}
\datecreate{2020-01-20}
\niveau{L1}
\module{Géométrie}
\chapitre{Géométrie de l'espace}
\sousChapitre{Distance}
\difficulte{3}
\theme{Distance}
\exo7qcmid{lille-1.geometrie-espace.029}

\contenu{
\texte{
Pour ces questions, l'espace est muni d'un repère orthonormé direct  $(O, \vec{i}, \vec{j}, \vec{k})$.
}

\begin{question}
\qclassification{141.05}
\qtags{motcle=distance}
Soit  $E$ l'ensemble des points situés à la même distance  des axes de coordonnées. Quelles sont les assertions vraies ?
\begin{answers}

       \bad{$E$ est une droite.}

       \good{$E$ est une réunion de droites.}

       \bad{$M(x,y,z) \in E \Leftrightarrow x=y=z$}

       \good{$M(x,y,z) \in E \Leftrightarrow |x|=|y|=|z|$}

\end{answers}
\begin{explanations}
$M(x,y,z) \in E \Leftrightarrow   \Vert \overrightarrow {OM} \wedge \vec{i} \Vert
 =  \Vert \overrightarrow {OM} \wedge \vec{j} \Vert =  \Vert \overrightarrow {OM} \wedge \vec{k} \Vert$.
\end{explanations}
\end{question}
}
